\documentclass[11pt]{article}
\usepackage{amsmath,amssymb,amsthm,mathtools,geometry,hyperref}
\geometry{margin=1in}
\newtheorem{theorem}{Theorem}
\newtheorem{proposition}{Proposition}
\newcommand{\C}{\mathbb C}
\newcommand{\Disc}{\operatorname{Disc}}
\newcommand{\sgn}{\operatorname{sgn}}
\newcommand{\Aut}{\operatorname{Aut}}
\title{Exact signed-root monodromy of the ES--Fable inverse cover}
\author{Independent programme calculation}
\date{20 September 2026}
\begin{document}
\maketitle

This calculation acts on the actual degree-eight inverse cover of
the original four-dimensional ES--Fable map, with the four marked
labels $\mathrm M,\mathrm N,\mathrm T,\mathrm E$ retained.
The polynomial and factor chart come from \cite{ES488}; the complete
global inverse and its chart boundary were proved in \cite{Global}.
All continuation paths, sign changes, group generators, and deck
transformations used here are constructed below. No identification
with a named Weyl group is assumed.

\section{The actual cover and its eight labelled states}
For $u=(u_0,u_2,u_3,u_4)$ put
\[
 H_u(U,V)=u_0U^4+U^3V+u_2U^2V^2+u_3UV^3+u_4V^4,
 \qquad h_u(x)=H_u(x,1),\qquad A=u_0.
\]
The original cubic coefficient remains one throughout. Let
\[
 \mathcal B=\{u:A\Disc(H_u)\ne0\},\qquad
 R=\C[u_0,u_2,u_3,u_4,(A\Disc H_u)^{-1}].
\]
Over this domain the original source is isomorphic to the finite
cover with coordinate algebra
\begin{equation}
 S=R[r,a]/(h_u(r),\ a^2h_u'(r)-1).\tag{SM1}
\end{equation}
In the unchanged original source coordinates its inverse is
\begin{align}
 y&=-r-i/a,\nonumber\\
 z&=A/a^3+2r/a+3i/a^2,\nonumber\\
 w&=7ir^2/a+(u_2-17r+Ar^2)/a^2-13i/a^3-2A/a^4.
 \tag{SM2}
\end{align}
These formulas are obtained by factoring
$H_u=(aU+bV)(cU^3+dU^2V+eUV^2+fV^3)$, setting $r=-b/a$,
and retaining the resultant equation
$C(-b,a)=a^2h_u'(r)=1$. Thus SM1 describes the actual inverse
states; its two signs are two distinct original source points.

For completeness, $R[r]/(h_u)$ is free of rank four because $A$
is a unit and $A^{-1}h_u$ is monic. The discriminant condition
and the Sylvester Bezout identity make $h_u'(r)$ a unit. The
second equation of SM1 is therefore the monic relation
$a^2=h_u'(r)^{-1}$, and $S$ is free of rank eight over $R$.
The Jacobian of the two equations with respect to $(r,a)$ has
determinant $2a h_u'(r)^2$, a unit. In a sufficiently small target
neighbourhood the four simple roots and their two nonzero square
roots therefore give eight separate holomorphic inverse charts.
This provides the local path-lifting property used below.

Both the base and the actual cover are path connected. The base
is the complement in $\C^4$ of the zero set of the nonzero
polynomial $A\Disc H_u$. The actual source is the complement
of the zero set of its polynomial pullback under the original map
$P$. This pullback is nonzero, since the eight states of any point
in $\mathcal B$ exist by SM1--2. To connect two points in either
complement, restrict its defining polynomial to the complex line
through them. Its value at the endpoints is nonzero, so the
restricted polynomial is not identically zero and has only
finitely many zeros. A path in the complex line parameter from
zero to one, with small detours around these finitely many points,
connects the endpoints within the complement. This proves the
asserted connectedness without an unproved topological
classification of the base.

\section{The exact upper bound from the retained product invariant}
At a base point order the four roots as $r_1,r_2,r_3,r_4$ and
choose one of the two actual values $a_j$ at each root. Denote
the eight states by $(j,+)$ and $(j,-)$, corresponding respectively
to $a_j$ and $-a_j$ in SM2. During continuation the four roots
remain distinct, and their scalings remain nonzero. Put
\[
 \Delta=\prod_{j<k}(r_k-r_j),\qquad
 \chi=A^2\Delta\prod_{j=1}^4a_j.
\]
The complete derivative product is
\[
 \prod_{j=1}^4h_u'(r_j)
 =A^4\prod_{j=1}^4\prod_{k\ne j}(r_j-r_k)
 =A^4\Delta^2.
\]
There are six pairs, so the sign from reversing each pair is
$(-1)^6=1$. Consequently
\begin{equation}
 \chi^2=1.\tag{SM3}
\end{equation}
The continuous value $\chi$ is therefore constant along every
ordered lift of a path. Its initial value may be $1$ or $-1$;
neither value has been discarded.

A loop in $\mathcal B$ induces a root permutation $\pi\in S_4$
and signs $\sigma_j\in\{1,-1\}$ such that
$r_j\mapsto r_{\pi(j)}$ and $a_j\mapsto\sigma_j a_{\pi(j)}$.
The negative state continues as the negative of the positive one,
by uniqueness in the equation $a^2h'(r)=1$. The exact action is
therefore
\begin{equation}
 (j,\epsilon)\longmapsto(\pi(j),\sigma_j\epsilon),
 \qquad \epsilon\in\{1,-1\}.\tag{SM4}
\end{equation}
The final Vandermonde is $\sgn(\pi)\Delta$. Applying the
constancy of SM3 gives the necessary relation
\begin{equation}
 \prod_{j=1}^4\sigma_j=\sgn(\pi).\tag{SM5}
\end{equation}
Thus the monodromy is contained in the explicitly defined group
\begin{equation}
 G=\{(\sigma,\pi)\in\{\pm1\}^4\rtimes S_4:
                       \prod_j\sigma_j=\sgn(\pi)\}.
 \tag{SM6}
\end{equation}
For each permutation exactly eight sign choices satisfy SM5, so
$|G|=8\cdot24=192$. As signed permutation matrices its elements
have determinant $(\prod_j\sigma_j)\sgn(\pi)=1$.

\section{Explicit local exchanges realizing the full group}
Use the real roots
\[
 (r_{\mathrm T},r_{\mathrm E},r_{\mathrm M},r_{\mathrm N})
 =(7,9,11,13),\qquad A=-1/40.
\]
Their sum is forty; hence the cubic coefficient of
$A\prod(x-r_j)$ is exactly one. The base target is
\begin{equation}
 u_*^{\rm real}=(-1/40,-59/4,95,-9009/40)\in\mathcal B.
 \tag{SM7}
\end{equation}
For the next calculation index the real order as
$s_1=7,s_2=9,s_3=11,s_4=13$; these indices refer respectively
to the unchanged labels $\mathrm T,\mathrm E,\mathrm M,\mathrm N$.
The four derivatives are
\[
 (h'(s_1),h'(s_2),h'(s_3),h'(s_4))
       =(6/5,-2/5,2/5,-6/5).
\]
Choose the exact base scalings
\begin{equation}
 (a_{\mathrm T},a_{\mathrm E},a_{\mathrm M},a_{\mathrm N})
 =\left(\sqrt{5/6},-i\sqrt{5/2},-\sqrt{5/2},i\sqrt{5/6}\right),
 \tag{SM8}
\end{equation}
where the real square roots displayed are positive. Each satisfies
the original equation in SM1.

For each adjacent pair $(s_j,s_{j+1})$, let
$c=(s_j+s_{j+1})/2$ and $d=1$. Keeping the other two roots fixed,
take the explicit path
\begin{equation}
 s_j(\theta)=c-d e^{i\theta},\qquad
 s_{j+1}(\theta)=c+d e^{i\theta},\qquad0\le\theta\le\pi.
 \tag{SM9}
\end{equation}
The pair sum remains $2c$, so the total root sum and the original
coefficient $A=-1/40$ are unchanged. The two moving roots remain
distinct. The closed disk of radius one about $c$ contains no
other root, so no collision with a fixed root occurs. At the
endpoint the pair is interchanged. The unordered roots and all
polynomial coefficients consequently return to the same target;
SM9 is a loop in the actual base $\mathcal B$.

We compute its signs without assuming a braid-group presentation.
For a fixed root $s_k$ outside the pair disk, the moving contribution
to $h'(s_k)$ is
\[
 (s_k-s_j(\theta))(s_k-s_{j+1}(\theta))
       =(s_k-c)^2-d^2e^{2i\theta}.
\]
Its values lie in the disk with positive real centre $(s_k-c)^2$
and strictly smaller radius $d^2$. This disk is in the right
half-plane and has a single continuous logarithm. After the loop
that logarithm returns to its initial value. Therefore neither
scaling at a fixed root changes sign.

For a moving root, its difference from the other moving root turns
by exactly $\pi$. Each of its two differences from an outside
root lies in a disk not containing zero, with both endpoint values
real and of the same sign. A logarithm on the corresponding left
or right half-plane shows that the net imaginary part of each of
those logarithm differences is zero. The argument change of the
complete moving derivative is thus exactly $\pi$. Its inverse
square root changes phase by $-\pi/2$, with its magnitude taking
the new root's required magnitude. By the explicit choices SM8,
the resulting permutation is the four-cycle
\begin{equation}
 \beta_j=(j+\quad(j+1)+\quad j-\quad(j+1)-).
 \tag{SM10}
\end{equation}
Equivalently, in the four retained Time labels the three generators
are
\begin{equation}
 \begin{split}
 \beta_1&=(\mathrm T+\ \mathrm E+\ \mathrm T-\ \mathrm E-),\\
 \beta_2&=(\mathrm E+\ \mathrm M+\ \mathrm E-\ \mathrm M-),\\
 \beta_3&=(\mathrm M+\ \mathrm N+\ \mathrm M-\ \mathrm N-).
 \end{split}\tag{SM11}
\end{equation}
Every sign and every fixed state has been determined by the path.

The square $\beta_j^2$ flips exactly the two adjacent root signs.
These three pair flips generate all even sign patterns: as vectors
over $\mathbb F_2$ they are
$(1,1,0,0)$, $(0,1,1,0)$ and $(0,0,1,1)$, which are independent
and span the dimension-three subspace with coordinate sum zero.
The projections of the three $\beta_j$ are the three adjacent
transpositions, and they generate all of $S_4$ by successively
moving entries to any prescribed positions. Thus the group
generated by these actual loops contains all eight elements in
the sign kernel and projects onto all 24 permutations. It has
at least $8\cdot24=192$ elements. The universal upper bound SM6
has exactly 192 elements, proving equality:
\begin{equation}
 \boxed{\operatorname{Mon}(P^{-1}(\mathcal B)/\mathcal B)=G.}
 \tag{SM12}
\end{equation}
This proof establishes the entire group without an unproved claim
that every base loop is a braid word.

\section{Transport from the literal original quartet}
Start with the original marking
\[
 (\rho_{\mathrm M},\rho_{\mathrm N},\rho_{\mathrm T},\rho_{\mathrm E})
 =\left(\tfrac12+\delta+i\gamma,\tfrac12+\delta-i\gamma,
        \tfrac12-\delta+i\gamma,\tfrac12-\delta-i\gamma\right),
 \qquad0<\delta<\tfrac12,\quad\gamma>2.
\]
Its target is $H_\rho=-\frac12\prod(U-\rho_jV)$, with the
original cubic coefficient one. The following two paths give a
fully explicit base transport to SM7. First, for $0\le t\le1$,
put
\begin{equation}
 \begin{array}{ll}
 r_{\mathrm M}(t)=\tfrac12+(1-t)(\delta+i\gamma)+t,
 &r_{\mathrm N}(t)=\tfrac12+(1-t)(\delta-i\gamma)+3t,\\
 r_{\mathrm T}(t)=\tfrac12+(1-t)(-\delta+i\gamma)-3t,
 &r_{\mathrm E}(t)=\tfrac12+(1-t)(-\delta-i\gamma)-t.
 \end{array}\tag{SM13}
\end{equation}
Their sum is always two, so $A=-1/2$ throughout this segment.
Roots with different original imaginary signs have different
imaginary parts for $t<1$. Roots with the same original imaginary
sign have real difference $2\delta(1-t)+4t$, which is positive.
At $t=1$ the four real roots are $(3/2,7/2,-5/2,-1/2)$ and
are distinct. Thus the entire segment remains in $\mathcal B$.

Next translate all four roots by $19t/2$, $0\le t\le1$, and
put $A(t)=-1/(2+38t)$. Their differences remain fixed, their sum
is $2+38t\ne0$, and the retained original cubic coefficient is
again one. At its endpoint the marked roots are
$(11,13,7,9)$, namely SM7 in the original Time order. This
second segment also lies entirely in $\mathcal B$.

The actual signs transport by the original equations, not by
discarding a branch. An exact formula along either differentiable
segment is
\begin{equation}
 a_j(t)=a_j(0)\exp\left[-\frac12\int_0^t
 \left(\frac{A'(s)}{A(s)}+
 \sum_{k\ne j}\frac{r_j'(s)-r_k'(s)}{r_j(s)-r_k(s)}\right)ds\right].
 \tag{SM14}
\end{equation}
The integrand has no singularity on these paths. Differentiating
shows that $a_j(t)^2A(t)\prod_{k\ne j}(r_j(t)-r_k(t))$ is
constant, and its initial value is one. Hence SM14 is precisely
the lifted inverse branch. Both initial signs lift separately.
Using SM8 as the signs at the real endpoint defines their exact
transported signs at the original quartet; any other initial
sign convention conjugates SM11 by a diagonal sign permutation.
It does not change SM12, because the determinant-one subgroup
is invariant under that conjugation.

Conjugating each loop SM9 by the explicitly constructed path
SM13 and its translation segment gives based loops at the literal
quartet with the same marked-root action. This transport concerns
the auxiliary quartic targets of the fixed nonlinear map. It does
not assert that intermediate roots are arithmetic zeta zeros or
that an original conductor period is moved along these paths.

\section{Why this group is not the Weyl group of type $D_4$}
With composition acting first by the rightmost generator, the
product $\beta_1\beta_2\beta_3$ is the eight-cycle
\begin{equation}
 (\mathrm T+\ \mathrm E+\ \mathrm M+\ \mathrm N+\
   \mathrm T-\ \mathrm E-\ \mathrm M-\ \mathrm N-).
 \tag{SM15}
\end{equation}
This follows by applying the three four-cycles SM11 to each
state in order. It has order eight.

For comparison, the reflections in the roots
$\{\pm e_j\pm e_k:j\ne k\}$ generate precisely the signed
permutations with an even total number of sign changes. Indeed
the reflection in $e_j-e_k$ swaps those coordinates, and the
reflection in $e_j+e_k$ swaps them with both signs changed.
Their compositions generate every permutation and every pair
sign flip, and each such reflection has even total sign product.
This proves the usual signed-permutation description of the
Weyl group $W(D_4)$ directly:
\[
 W(D_4)=\{(\sigma,\pi):\prod_j\sigma_j=1\}.
\]
This group has no element of order eight. An underlying
permutation of four coordinates can have cycle lengths only
$1,2,3,4$. A signed cycle of length $m$ has order $m$ if its
sign product is positive and $2m$ if its sign product is
negative, as follows by applying the cycle $m$ times to each
coordinate. With no length-four cycle, the available orders
combine to divisors of four or six. With a length-four cycle,
the even total sign condition makes its sign product positive,
so its order is four. Thus order eight never occurs in $W(D_4)$,
whereas SM15 occurs in $G$. The groups are therefore not even
abstractly isomorphic, despite their equal order 192.

\section{Exact deck groups}
Let $\tau$ denote simultaneous reversal of all four root signs,
$\tau(j,\epsilon)=(j,-\epsilon)$. The centralizer of $G$ in
the full permutation group of the eight inverse states is
\begin{equation}
 C_{S_8}(G)=\{1,\tau\}.\tag{SM16}
\end{equation}
Here is a complete proof. Fix one state, say $(1,+)$. Its
stabilizer in $G$ also fixes $(1,-)$. No state at another root
is fixed by that stabilizer: a pair sign flip involving that
root and one of the other two roots different from root 1
belongs to the stabilizer and moves that state. Hence the
stabilizer has exactly the two fixed states $(1,+),(1,-)$.
Any permutation $c$ commuting with $G$ must send $(1,+)$ to
one of these fixed states. The group $G$ is transitive on all
eight states (permutations move root labels and pair flips
change either chosen sign), so the image of this one state
determines $c$ everywhere through $c(gx)=g c(x)$. The first
choice gives the identity and the second gives $\tau$.
Both commute with every signed permutation, proving SM16.

A deck transformation of the actual connected cover restricts
on the base fibre to a permutation commuting with monodromy:
apply the deck transformation to a lifted loop and use uniqueness
of path lifting in the local charts of SM1. Its restriction is
injective as a function of the deck transformation. Indeed a
transformation fixing one fibre point fixes the lift of every
path starting there; connectedness, or the explicit path-lifting
description, reaches the entire cover. Thus SM16 is an upper
bound of two for the actual deck group.

Both transformations are realized. The nontrivial one is the
original rational factor-sign involution, in unchanged source
coordinates,
\begin{equation}
 \Sigma(a,y,z,w)=\left(-a,\ y+\frac{2i}{a},\
 \frac{6i}{a^2}-z,\
 w-\frac{14iy^2}{a}+\frac{28y}{a^2}+\frac{40i}{a^3}\right).
 \tag{SM17}
\end{equation}
To check it directly in the inverse model, keep the same $r,u$
in SM2 and replace $a$ by $-a$. Its new $y$ is $y+2i/a$ and
its new $z$ is $6i/a^2-z$. Expanding the two expressions for
$w$ gives exactly SM17. It follows from this same operation
that $\Sigma^2=1$ and $P\Sigma=P$ on $a\ne0$. Since $u_0=ac$
is nonzero on $\mathcal B$, every point of its preimage has
$a\ne0$. SM17 is regular on the entire cover and realizes
$\tau$. The exact cover deck group is consequently
\begin{equation}
 \Aut_{\mathcal B}(P^{-1}(\mathcal B))=\{1,\Sigma\}\cong C_2.
 \tag{SM18}
\end{equation}

The same two elements are all rational birational deck
transformations of the original polynomial map. To justify the
passage from rational maps, write $K_0=\C(u_0,u_2,u_3,u_4)$
and $L_0=\C(a,y,z,w)$ with the embedding defined by $P$.
SM1 makes $[L_0:K_0]=8$. The actual ring $S$ is also the
localization
\[
 \C[a,y,z,w,(P_0\Disc(H_{P(q)}))^{-1}],
\]
so it is integrally closed: a localization of a polynomial
unique-factorization ring has this property, since a reduced
fraction integral over it has denominator dividing a power of
its relatively prime numerator and therefore has unit denominator.
Since $S$ is finite integral over $R$, it is the full integral
closure of $R$ in $L_0$. In fact any element of $L_0$ integral
over $R$ is integral over $S$ by the same monic polynomial,
and integrally closedness puts it in $S$.
Every $K_0$-automorphism of $L_0$ preserves this integral closure,
and hence extends to a regular deck transformation on the actual
cover. SM18 therefore bounds it by those same two elements;
SM17 realizes both. Thus
\begin{equation}
 \Aut_{K_0}(L_0)=\{1,\Sigma\}.\tag{SM19}
\end{equation}
This statement is about birational automorphisms fixing the
target field, not arbitrary rational maps of the source.

Finally the only global polynomial deck automorphism $D$ of $\C^4$
satisfying $P\circ D=P$ is the identity. Such an automorphism would
restrict to SM19. The nonidentity element cannot extend
polynomially across $a=0$, because its second coordinate is
$(ay+2i)/a$, whose numerator is not divisible by $a$.
Equality on the dense open set $a\ne0$ forces equality as
rational functions, so there is no alternative polynomial
extension. This proves the assertion for the whole original
source, while retaining the nontrivial regular deck
transformation on the precise eight-sheet domain. This does not
classify polynomial automorphisms satisfying the different equation
$P\circ D=D\circ P$; that centralizer is not asserted to be trivial.

\section{Receiving the group in the fixed original conductor target}
For the fixed original receiving isomorphism
$\Psi:\C^4\to W=\cP_3$, with all original moments and Gamma
norms retained, set $P_W=\Psi P\Psi^{-1}$. Its eight-sheet
base is exactly $\Psi(\mathcal B)$. Sending a target path
$u(t)$ to $\Psi u(t)$ and each source state to its image under
$\Psi$ gives a bijection of all path lifts, with inverse
$\Psi^{-1}$. Thus its monodromy is exactly SM12 on the eight
transported marked states, and its nontrivial rational deck
transformation is exactly $\Psi\Sigma\Psi^{-1}$. The cover
deck group is $C_2$ and its global polynomial deck group is
trivial by the same conjugacy.

These are statements about the constructed nonlinear cover on
the original target. They do not enlarge the symmetry group
of a fixed Segre pairing or assert that the original conductor
itself is a nonlinear map. The original linear conductor
remains the exact invertible intertwiner proved in
\cite{Bridge}. All four Time labels and both signs over each
are present in the monodromy action just calculated.

\begin{thebibliography}{9}
\bibitem{ES488} The Clankers, \emph{Erd\H{o}s--Straus project reader},
canonical 488-page edition, deposited source version 69,
\texttt{erdos\_straus\_project\_reader.tex}, theorem
\texttt{explicit-four-time-Jacobian-collision}, source lines
19484--19588. The source attribution identifies the original
four-dimensional map and its marked coordinates; the monodromy
calculation in this note is not retrospectively attributed to it.
\bibitem{Global} Programme calculation,
\emph{Complete fibres, the eight-sheet domain, and the nonproper-value
locus of the ES--Fable map},
\texttt{GLOBAL\_FIBRE\_AND\_NONPROPER\_LOCUS.tex}, GF1--28,
20 September 2026. SM1--2 are its exact inverse cover and original
coordinate formulas; the group and deck proofs here are new.
\bibitem{Bridge} Programme calculation,
\emph{The four ES--Fable locales and the original weighted conductor},
\texttt{FABLE\_TO\_ORIGINAL\_CONDUCTOR.tex}, FC26--35 and FC50--55,
20 September 2026. The original WCF operator and its human
orthogonal-polynomial provenance remain as cited there.
\end{thebibliography}
\end{document}
